\booksection{Appendice al Capitolo 11: Conoscenze di tecnica professionale dell'immagine}

\chapterintro{Per comprendere davvero i segreti della fotografia e dei video sul telefono, devi conoscere alcune nozioni professionali di tecnica dell'immagine. Queste conoscenze non solo ti aiutano a usare meglio la fotocamera del telefono, ma ti rendono anche più consapevole nella scelta dei dispositivi.}

\booksubsection{Risoluzione, frame rate e bitrate: i tre elementi della qualità video}

\subsubsection*{Risoluzione (Resolution):}

La risoluzione determina la nitidezza dell'immagine, espressa in numero di pixel.

\begin{itemize}
  \item \textbf{720P (HD)}: 1280×720 pixel, nitidezza entry-level.
  \item \textbf{1080P (Full HD)}: 1920×1080 pixel, la risoluzione più diffusa attualmente.
  \item \textbf{2K (QHD)}: 2560×1440 pixel, più nitido di 1080P.
  \item \textbf{4K (UHD)}: 3840×2160 pixel, ultra HD, ricco di dettagli.
  \item \textbf{8K}: 7680×4320 pixel, standard futuro, supporto dispositivo attualmente limitato.
\end{itemize}

Più alta è la risoluzione, più nitida è l'immagine, ma il file è più grande e le esigenze di archiviazione ed elaborazione aumentano.

\subsubsection*{Frame rate (Frame Rate):}

Il frame rate determina la fluidità dell'immagine, espresso in frame per secondo (fps).

\begin{itemize}
  \item \textbf{24fps}: frame rate standard del cinema, con una peculiare "sensazione cinematografica".
  \item \textbf{30fps}: comune per la TV e i video web, fluido e naturale.
  \item \textbf{60fps}: frame rate alto, adatto a scene sportive e registrazione di gioco, molto fluido.
  \item \textbf{120fps/240fps}: frame rate altissimo, utilizzabile per video al rallentatore.
\end{itemize}

Più alto è il frame rate, più fluida è l'immagine e migliore è l'effetto slow motion, ma il file è più grande.

\subsubsection*{Bitrate:}

Il bitrate determina la densità di dati del video, espresso in dati al secondo (Mbps).

\begin{itemize}
  \item Più alto è il bitrate, migliore è la qualità, maggiori i dettagli conservati, ma il file è più grande.
  \item Un bitrate troppo basso provoca nell'immagine mosaicature, blocchi di colore e altri artefatti di compressione.
  \item Le riprese video professionali richiedono in genere un bitrate alto per conservare più spazio in post-produzione.
\end{itemize}

\booksubsection{Formati di codifica video: H.264, H.265 e ProRes}

\subsubsection*{H.264 (AVC):}

\begin{itemize}
  \item \textbf{Caratteristiche}: il formato di codifica più diffuso, compatibilità eccellente.
  \item \textbf{Vantaggi}: quasi tutti i dispositivi possono riprodurlo, efficienza di compressione moderata.
  \item \textbf{Svantaggi}: rispetto ai nuovi formati, efficienza di compressione più bassa, file più grandi.
  \item \textbf{Scenari applicativi}: condivisione di video online, riprese quotidiane degli utenti comuni.
\end{itemize}

\subsubsection*{H.265 (HEVC):}

\begin{itemize}
  \item \textbf{Caratteristiche}: la versione aggiornata di H.264, efficienza di compressione migliorata di circa il 50\%.
  \item \textbf{Vantaggi}: a parità di qualità il file è più piccolo, supporta risoluzioni e frame rate più alti.
  \item \textbf{Svantaggi}: la codifica e la decodifica richiedono processori più potenti, alcuni dispositivi datati non sono compatibili.
  \item \textbf{Scenari applicativi}: riprese video 4K, scenari con spazio di archiviazione limitato.
\end{itemize}

\subsubsection*{Apple ProRes:}

\begin{itemize}
  \item \textbf{Caratteristiche}: codifica video professionale sviluppata da Apple, compressione quasi lossless.
  \item \textbf{Vantaggi}: qualità dell'immagine altissima, ampio spazio di post-produzione, colori accurati.
  \item \textbf{Svantaggi}: file enormi, richiede molto spazio di archiviazione, l'elaborazione richiede attrezzature professionali.
  \item \textbf{Scenari applicativi}: produzione video professionale, riprese di livello cinematografico.
\end{itemize}

\subsubsection*{ProRes RAW:}

\begin{itemize}
  \item \textbf{Caratteristiche}: combina i vantaggi di ProRes e del formato RAW.
  \item \textbf{Cos'è il RAW}: dati originali del sensore non elaborati, conserva la massima quantità di informazioni.
  \item \textbf{Vantaggi}: spazio di regolazione in post-produzione enorme, si possono modificare parametri come esposizione e bilanciamento del bianco senza perdita di qualità.
  \item \textbf{Svantaggi}: file giganteschi, richiedono software di post-produzione professionali per l'elaborazione.
  \item \textbf{Scenari applicativi}: produzioni professionali di alto livello. iPhone 15 Pro serie e successivi supportano la registrazione video Apple ProRes (fino a 4K 60fps), iPhone 17 Pro supporta ulteriormente la registrazione ProRes RAW direttamente in macchina.
\end{itemize}

\booksubsection{HDR: tecnologia High Dynamic Range}

La tecnologia HDR (High Dynamic Range) consente di conservare simultaneamente i dettagli nelle parti più chiare e più scure dell'immagine.

\begin{itemize}
  \item \textbf{HDR10}: lo standard HDR più basilare, buona compatibilità.
  \item \textbf{HDR10+}: standard promosso da Samsung, supporta metadati dinamici.
  \item \textbf{Dolby Vision}: standard HDR di punta sviluppato da Dolby, qualità video migliore, ma richiede licenza.
  \item \textbf{HLG}: standard HDR broadcast, compatibile con display comuni.
\end{itemize}

Quando si riprendono video HDR, serve un display che supporti HDR per vedere l'effetto completo.

\booksubsection{Spazio colore e profondità di colore}

\subsubsection*{Spazio colore:}

\begin{itemize}
  \item \textbf{sRGB}: spazio colore standard, la migliore compatibilità.
  \item \textbf{Display P3}: standard wide gamut promosso da Apple, colori più ricchi.
  \item \textbf{Rec.2020}: spazio colore standard per video HDR, la gamma più ampia.
\end{itemize}

\subsubsection*{Profondità di colore:}

\begin{itemize}
  \item \textbf{8-bit}: 16,7 milioni di colori, standard per i video comuni.
  \item \textbf{10-bit}: 1,07 miliardi di colori, standard per video professionali, transizioni cromatiche più uniformi.
  \item \textbf{12-bit}: 68 miliardi di colori, standard professionale di fascia alta.
\end{itemize}

\begin{tcbinfo}{Consigli pratici}
\begin{itemize}
  \item Riprese quotidiane: 1080P/4K + 30fps + codifica H.265, equilibra qualità e dimensione del file.
  \item Riprese sportive: 1080P/4K + 60fps, garantisce fluidità.
  \item Creazione professionale: 4K + ProRes (se supportato), conserva il massimo spazio in post-produzione.
  \item Consiglio di archiviazione: i video ad alto bitrate occupano molto spazio, conviene fare backup periodici su computer o cloud.
\end{itemize}

\end{tcbinfo}

\begin{tcbpractice}{Esplora: le impostazioni video del tuo telefono}
\begin{enumerate}
  \item Apri l'app Fotocamera, passa alla modalità video.
  \item Controlla le opzioni di risoluzione e frame rate, prova combinazioni diverse.
  \item Se supportato, verifica se ci sono opzioni avanzate come HDR, ProRes.
  \item Registra alcuni spezzoni con impostazioni diverse e confronta qualità e dimensione del file.
\end{enumerate}

\end{tcbpractice}

\booksubsection{Formato RAW vs HEIF/JPEG}

\subsubsection*{Che cos'è il formato RAW:}

Il formato RAW è il “dato originale” registrato dal sensore della fotocamera - un'informazione pura non compressa ed elaborata. È come un ingrediente crudo, che richiede una “cottura” in post-produzione per essere ultimato. Poiché il telefono ha memoria e potenza di calcolo limitate, per impostazione predefinita scatta JPEG o HEIF “già cotti”; il RAW va attivato appositamente.

\subsubsection*{Confronto RAW vs HEIF vs JPEG:}

\begin{itemize}
  \item \textbf{JPEG}: il formato più vecchio, la migliore compatibilità (qualsiasi dispositivo può aprirlo). Profondità colore 8-bit, compressione lossy, il file più piccolo (2-5MB a foto). Spazio di post-produzione limitato, forzando l'illuminazione delle ombre compaiono facilmente blocchi di colore.
  \item \textbf{HEIF/HEIC}: nuovo formato proposto da Apple, a parità di qualità il file è la metà di JPEG (1-3MB a foto). Supporta profondità colore 10-bit e canale alfa. Lo svantaggio è che alcuni dispositivi datati e Windows richiedono plug-in per aprirlo.
  \item \textbf{RAW}: dato originale non compresso, il file più grande (20-100MB a foto). Profondità colore di solito 12-16bit, massimo spazio di post-produzione. Lo svantaggio è che occorre software professionale (Lightroom, Snapseed, ecc.) per aprirlo ed elaborarlo, non può essere condiviso direttamente.
\end{itemize}

\subsubsection*{Perché i fotografi professionisti scattano in RAW:}

\begin{itemize}
  \item \textbf{Bilanciamento del bianco regolabile in post senza perdite}: se JPEG ha il bilanciamento del bianco errato (ad esempio tendente al giallo sotto lampada a incandescenza), la correzione in post-produzione fa perdere qualità; con RAW si può reimpostare la temperatura colore, senza alcuna perdita.
  \item \textbf{Ampia tolleranza di esposizione}: in JPEG le ombre sottoesposte di 2 stop sono praticamente nere, schiarendole compare solo rumore; il RAW può "salvare" dettagli utilizzabili dalle ombre sottoesposte di 3-4 stop. Lo stesso vale per le sovraesposizioni.
  \item \textbf{Dettagli più ricchi in ombre e alte luci}: la gamma dinamica del RAW è di gran lunga superiore a JPEG, molto maggiore è il margine per evitare bruciati nelle luci e neri chiusi nelle ombre.
  \item \textbf{Nitidezza, riduzione rumore, correzione obiettivo tutti controllabili in post}: l'elaborazione interna della fotocamera è "una volta sola", mentre con RAW si può intervenire più volte, sempre partendo dai dati originali.
\end{itemize}

\subsubsection*{Come attivare il RAW sul telefono:}

\begin{itemize}
  \item \textbf{iPhone ProRAW}: supportato solo da iPhone 12 Pro e modelli Pro superiori. Apri "Impostazioni Fotocamera Formato" e attiva "Apple ProRAW". In alto nell'interfaccia della fotocamera comparirà un pulsante "RAW", attivalo. Circa 25MB per foto, consiglio di attivarlo solo quando serve.
  \item \textbf{Android Motion Raw / Ultra RAW}: ogni marchio lo chiama in modo diverso. Su Xiaomi/OPPO/vivo di punta trovi l'interruttore RAW in "Modalità professionale" o "Altro" della fotocamera; su Huawei attiva RAW in "Modalità professionale"; la serie Samsung Galaxy S ha l'opzione RAW in modalità "Pro". Le dimensioni del file sono in genere 20-50MB.
  \item \textbf{Consiglio generale}: per gli scatti quotidiani non serve attivare RAW, il file è troppo grande e occupa spazio; attivalo solo in scenari importanti (paesaggi, ritratti, con necessità di post-produzione).
\end{itemize}

\booksubsection{Principi della fotografia computazionale}

La fotografia moderna del telefono non può più fare a meno della “fotografia computazionale” - ovvero affidarsi agli algoritmi e non solo a lenti e sensori per generare le foto. Capire questi principi chiave ti fa comprendere meglio cosa fa il telefono nell'istante in cui premi l'otturatore.

\subsubsection*{Fusione multi-frame:}

Quando premi l'otturatore sul telefono, in realtà scatta più foto. L'algoritmo allinea e fonde i frame pesandoli, per ottenere uno scatto finale di alta qualità.

\begin{itemize}
  \item \textbf{Modalità notturna}: in 1-3 secondi scatta una raffica di foto sottoesposte, le allinea e le somma - il rumore si annulla a vicenda e la luminosità si accumula. Il risultato è che anche in bassa luce si ottengono immagini pulite.
  \item \textbf{HDR (High Dynamic Range)}: scatta separatamente tre (o più) foto sottoesposte, normali e sovraesposte, e le fonde in un'unica foto con dettagli sia nelle luci che nelle ombre. Particolarmente efficace in controluce.
  \item \textbf{Deep Fusion (Apple) / Photonic Engine}: per scene a luce medio-bassa, analizza e fonde multi-frame pixel per pixel, con focus sul miglioramento di texture e dettagli; dettagli come maglioni e capelli risulteranno visibilmente più nitidi.
\end{itemize}

\subsubsection*{Segmentazione semantica:}

L'AI non elabora l'immagine in blocco, ma prima “capisce” cosa c'è nell'inquadratura - cielo, volti, prato, edifici, vestiti - e poi applica strategie diverse a ciascuna zona.

\begin{itemize}
  \item \textbf{Cielo}: aumenta il blu e abbassa la luminosità separatamente, per evitare sovraesposizione.
  \item \textbf{Volto}: riconosce le zone di tonalità della pelle, le schiarisce separatamente e uniforma la pelle, senza toccare i lineamenti e i dettagli dei capelli.
  \item \textbf{Prato/foglie}: aumenta la saturazione del verde, incrementa la nitidezza dei dettagli.
  \item \textbf{Sfondo sfocato}: in modalità ritratto, l'AI riconosce il contorno del soggetto, sfoca lo sfondo e mantiene il soggetto nitido.
\end{itemize}

\subsubsection*{Applicazioni del deep learning:}

\begin{itemize}
  \item \textbf{Riconoscimento scenico AI}: la fotocamera riconosce automaticamente scene come "gastronomia", "animali", "notte", "documento" e ottimizza i parametri di conseguenza. Ad esempio la modalità cibo aumenta le tonalità calde e la saturazione per rendere il cibo più invitante.
  \item \textbf{Sfocatura ritratto AI}: le lenti del telefono fisicamente non possono ottenere una profondità di campo ridotta come le reflex; si affida a modelli di deep learning per prevedere i bordi del soggetto e sfocare lo sfondo. Gli algoritmi più recenti sanno gestire correttamente capelli, bordi di vetri e altri dettagli.
  \item \textbf{Cancellazione AI / Espansione AI}: selezionando persone o oggetti nell'inquadratura, l'AI "ricostruisce" la parte nascosta in base al contesto. È un'applicazione dell'AI generativa nel campo dell'immagine.
\end{itemize}

\subsubsection*{Fotografia computazionale vs fotografia ottica:}

\begin{itemize}
  \item \textbf{Vantaggi della fotografia computazionale}: anche in volume ridotto si ottiene buona qualità, in scene estreme come bassa luce e controluce si riesce a scattare, gli utenti comuni senza nozioni di fotografia possono scattare foto "venute bene".
  \item \textbf{Svantaggi della fotografia computazionale}: a volte l'immagine è "sovraelaborata" - eccessiva nitidezza, colori alterati, colore del cielo innaturale; la fusione multi-frame lascia doppi contorni quando il soggetto si muove; i dettagli sono "calcolati" e non realmente registrati.
  \item \textbf{Vantaggi della fotografia ottica}: sensore grande + lente grandangolare catturano la luce in modo reale, qualità naturale, ampio spazio in post; non compaiono artefatti dovuti a "fantasie" algoritmiche.
  \item \textbf{Tendenza}: i telefoni di punta moderni combinano "ottica + calcolo" - usano sensori grandi per ottenere informazioni reali e algoritmi per ottimizzare la resa. Le due cose non sono in opposizione.
\end{itemize}

\booksubsection{Tipi di lente e lunghezze focali}

I telefoni di punta moderni sono in genere dotati di tre o quattro lenti, ciascuna corrispondente a una lunghezza focale diversa. Comprendere le lunghezze focali ti permette di scegliere la lente giusta in scenari diversi.

\subsubsection*{Lenti comuni e lunghezze focali equivalenti:}

\begin{itemize}
  \item \textbf{Grandangolo (13-24mm equivalente)}: il campo più ampio, adatto a paesaggi, architettura, spazi interni. Riesce a riprendere grandi scene davanti a te, ma i bordi presentano distorsione (linee dritte diventano curve) e nei ritratti il viso viene stirato. Adatto a "grandi scene", non a "primi piani".
  \item \textbf{Lente principale (24-28mm equivalente)}: la lente più usata quotidianamente, qualità migliore e apertura massima. Bilancia campo e qualità, adatta sia a ritratti sia a paesaggi. La "1x" della maggior parte dei telefoni è la lente principale.
  \item \textbf{Teleobiettivo (70-200mm equivalente)}: per soggetti lontani e ritratti in primo piano. Il tele comprime la sensazione di spazio (montagne lontane e persona vicina sembrano più vicine), nei ritratti si ottengono proporzioni del viso più naturali e sfondo sfocato.
  \item \textbf{Teleobiettivo periscopico (200mm+ equivalente)}: tramite un design a percorso ottico riflesso, realizza focali più lunghe nello spessore del telefono. Riesce a riprendere dettagli nitidi da molto lontano, adatto a concerti, eventi sportivi, fauna selvatica. Solo su modelli di punta.
\end{itemize}

\subsubsection*{Compromesso tra dimensione del sensore e dimensione dei pixel:}

Il sensore è il “fondo” e determina quanta luce può ricevere. Nel settore si dice “un fondo più grande schiaccia l'avversario” - a parità di numero di pixel, più grande è il sensore, più grandi sono i singoli pixel, maggiore è la luce raccolta, migliore è la qualità.

\begin{itemize}
  \item \textbf{Sensore grande}: molta luce in entrata, buona qualità notturna, ampia gamma dinamica. Lo svantaggio è che il telefono deve essere spesso più spesso (il modulo fotocamera sporgente), costo elevato.
  \item \textbf{Piccolo sensore ad alto numero di pixel}: di giorno, con buona luce, alti pixel permettono foto più nitide, agevoli per il ritaglio in post. Ma in bassa luce ciascun pixel riceve poca luce, con molto rumore.
  \item \textbf{Scelte dei flagship mainstream}: la fotocamera principale è in genere da 1/1,3 pollici a 1 pollice (alcuni modelli hanno già adottato sensori da 1 pollice), con circa 50 megapixel - questo è il punto di equilibrio tra qualità e volume.
\end{itemize}

\subsubsection*{Tecnologia di fusione pixel (Pixel Binning):}

Per bilanciare alti pixel di giorno e pixel grandi di notte, i produttori usano la “fusione pixel” - combina più pixel piccoli in un pixel grande in uscita.

\begin{itemize}
  \item \textbf{Pixel quad-bayer (4-in-1)}: sensore da 50 megapixel, foto in uscita a 12,5 megapixel. Ciascun pixel in uscita è sintetizzato da 4 pixel adiacenti, la luce in ingresso aumenta di 4 volte, qualità più pulita.
  \item \textbf{Pixel nonacell (9-in-1)}: per sensori da oltre 100 megapixel, uscita a circa 11-12 megapixel. Stesso principio, luce in ingresso maggiore.
  \item \textbf{Effetto pratico}: in modalità predefinita, si ottiene una foto a basso numero di pixel dopo la fusione - buona qualità e file piccolo. Quando servono dettagli estremi (documenti, soggetti lontani), si può passare manualmente in "modalità alto numero di pixel" per l'uscita a piena risoluzione.
\end{itemize}

\booksubsection{Modalità di ripresa speciali}

Oltre alle foto quotidiane, il telefono ha alcune modalità speciali; usate bene, permettono di ottenere immagini altrimenti impossibili.

\subsubsection*{Modalità panoramica:}

\begin{itemize}
  \item Durante la ripresa segui la freccia e muovi il telefono lentamente e in modo stabile in orizzontale, non veloce - troppo veloce causa errori di giunzione.
  \item Mantieni il telefono stabile in verticale, spostati solo in orizzontale. Se oscilli in alto-basso, l'orizzonte diventerà ondulato.
  \item Con persone o veicoli in movimento usala con cautela - nella giunzione possono comparire "mezzi corpi" o "fantasmi".
  \item La modalità panoramica può essere usata anche in verticale - ripresa verticale con giunzione orizzontale, adatta a edifici alti o cascate.
\end{itemize}

\subsubsection*{Slow motion e time-lapse:}

\begin{itemize}
  \item \textbf{Slow motion}: ripreso con frame rate alto a 120fps o 240fps, in playback diventa al ralenti. Adatto a gocce d'acqua, istanti sportivi, salti di animali. La luce deve essere sufficiente, in slow motion il tempo di esposizione è breve, in bassa luce la qualità scadente.
  \item \textbf{Time-lapse}: scatta foto a intervalli e le unisce in un video accelerato. Adatto a traffico, movimento delle nuvole, tramonti, fiori che si aprono. La durata di ripresa varia da pochi minuti a ore, fissa il telefono su un treppiede, a mano esce sfocato.
  \item \textbf{Piccolo trucco}: per time-lapse delle nuvole, posiziona il telefono e riprendi per 15-30 minuti, otterrai alla fine un video di circa 10 secondi, con un effetto d'impatto.
\end{itemize}

\subsubsection*{Modalità stelle/astronomica:}

\begin{itemize}
  \item Occorre allontanarsi dall'inquinamento luminoso cittadino, trovare una zona periferica o montana senza interferenze di lampioni. In città non si riesce a riprendere la Via Lattea.
  \item È indispensabile un treppiede, a mano esce sfocato. Il telefono deve restare stabile per alcuni secondi o decine di secondi.
  \item Il cielo deve essere sereno e senza nubi, la luna non troppo grande (la luna piena copre la luce delle stelle). Le notti buie di inizio o fine mese lunare sono le migliori.
  \item La "modalità notturna" di iPhone riconosce automaticamente il cielo stellato e prolunga il tempo di esposizione (fino a 30 secondi); la "modalità stelle" o "modalità astronomica" dei flagship Android ha effetti simili. Pazienta, dopo aver premuto l'otturatore non muovere il telefono.
\end{itemize}

\subsubsection*{Modalità macro:}

\begin{itemize}
  \item Per riprendere piccoli soggetti come fiori, insetti, gioielli, dettagli di testo. Può mettere a fuoco a 1-2 cm di distanza.
  \item Alcuni telefoni hanno una lente macro dedicata (spesso una lente piccola da 2 o 5 megapixel), con qualità media; i flagship usano la lente grandangolare anche per la macro, con qualità migliore.
  \item Riprendendo da molto vicino, la lente è molto vicina al soggetto e copre la luce; se necessario integra con lampada o torcia.
  \item La distanza di messa a fuoco è minuscola, la profondità di campo molto ridotta, basta che il soggetto si sposti di pochi millimetri per perdere la messa a fuoco. Mano ferma durante lo scatto, meglio se il soggetto resta immobile.
\end{itemize}

\begin{tcbinfo}{Consigli pratici per la fotografia statica}
\begin{itemize}
  \item Scatti quotidiani: HEIF/JPEG predefinito, file piccolo e qualità sufficiente.
  \item Scenari importanti (paesaggi, ritratti, post-produzione): attiva RAW, conserva il massimo spazio in post.
  \item Scegli la lente giusta: grandangolo per paesaggi, tele per ritratti, principale per uso quotidiano.
  \item Notte e stelle: fissa il telefono, attiva modalità notturna o astronomica.
  \item Telefoni con sensore grande (es. 1 pollice) hanno una qualità massima più alta, ma le abitudini operative restano determinanti.
\end{itemize}

\end{tcbinfo}

\begin{tcbpractice}{Prova: scopri le capacità della tua fotocamera}
\begin{enumerate}
  \item Verifica se il tuo telefono supporta lo scatto RAW (Impostazioni Fotocamera Formato), attivalo e scatta una foto, confronta le dimensioni del file JPEG e RAW della stessa scena.
  \item Passa alle tre focali grandangolo, principale e tele, scatta separatamente lo stesso soggetto lontano, apprezzando differenze di campo e qualità.
  \item Trova una scena con modalità notturna o stelle, fissa il telefono e scatta una foto notturna.
  \item Prova a scattare in modalità panoramica orizzontale una grande scena, facendo attenzione a muovere il telefono in modo stabile.
\end{enumerate}

\end{tcbpractice}
